% ══════════════════════════════════════════════════════════════════════════════
% 台積電獎學金申請文件之共用設定（字型與中文標籤）。
% 各份文件以 \documentclass[...,fontset=none]{ctexart} 開頭，再 \input 本檔：
%   % ══════════════════════════════════════════════════════════════════════════════
% 台積電獎學金申請文件之共用設定（字型與中文標籤）。
% 各份文件以 \documentclass[...,fontset=none]{ctexart} 開頭，再 \input 本檔：
%   % ══════════════════════════════════════════════════════════════════════════════
% 台積電獎學金申請文件之共用設定（字型與中文標籤）。
% 各份文件以 \documentclass[...,fontset=none]{ctexart} 開頭，再 \input 本檔：
%   % ══════════════════════════════════════════════════════════════════════════════
% 台積電獎學金申請文件之共用設定（字型與中文標籤）。
% 各份文件以 \documentclass[...,fontset=none]{ctexart} 開頭，再 \input 本檔：
%   \input{../common/preamble.tex}
% 版面（geometry、行距、頁碼）由各份文件自行設定。
% ══════════════════════════════════════════════════════════════════════════════

% ── JFM=zhtw (common/jfm-zhtw.lua): ctex's default JFM assumes simplified-
% Chinese glyphs (punctuation at the left of the em box), but the Traditional-
% Chinese fonts below centre ，：；。 etc., so compressed punctuation overlaps the
% next character. zhtw keeps those marks full width.
% ── CJK font: DFKai-SB (標楷體) on Windows, falling back to AR PL KaitiM Big5
% when DFKai-SB isn't installed (same setup as the midterm report).
\IfFontExistsTF{DFKai-SB}{%
  \setCJKmainfont{DFKai-SB}[JFM=zhtw]
  \setCJKsansfont{DFKai-SB}[JFM=zhtw]
  \setCJKmonofont{DFKai-SB}[JFM=zhtw]
}{%
  \setCJKmainfont{AR PL KaitiM Big5}[JFM=zhtw]
  \setCJKsansfont{AR PL KaitiM Big5}[JFM=zhtw]
  \setCJKmonofont{AR PL KaitiM Big5}[JFM=zhtw]
}

% ── Traditional-Chinese wording overrides ──────────────────────────────────────
\renewcommand{\figurename}{圖}
\renewcommand{\tablename}{表}
\renewcommand{\refname}{參考文獻}

% 版面（geometry、行距、頁碼）由各份文件自行設定。
% ══════════════════════════════════════════════════════════════════════════════

% ── JFM=zhtw (common/jfm-zhtw.lua): ctex's default JFM assumes simplified-
% Chinese glyphs (punctuation at the left of the em box), but the Traditional-
% Chinese fonts below centre ，：；。 etc., so compressed punctuation overlaps the
% next character. zhtw keeps those marks full width.
% ── CJK font: DFKai-SB (標楷體) on Windows, falling back to AR PL KaitiM Big5
% when DFKai-SB isn't installed (same setup as the midterm report).
\IfFontExistsTF{DFKai-SB}{%
  \setCJKmainfont{DFKai-SB}[JFM=zhtw]
  \setCJKsansfont{DFKai-SB}[JFM=zhtw]
  \setCJKmonofont{DFKai-SB}[JFM=zhtw]
}{%
  \setCJKmainfont{AR PL KaitiM Big5}[JFM=zhtw]
  \setCJKsansfont{AR PL KaitiM Big5}[JFM=zhtw]
  \setCJKmonofont{AR PL KaitiM Big5}[JFM=zhtw]
}

% ── Traditional-Chinese wording overrides ──────────────────────────────────────
\renewcommand{\figurename}{圖}
\renewcommand{\tablename}{表}
\renewcommand{\refname}{參考文獻}

% 版面（geometry、行距、頁碼）由各份文件自行設定。
% ══════════════════════════════════════════════════════════════════════════════

% ── JFM=zhtw (common/jfm-zhtw.lua): ctex's default JFM assumes simplified-
% Chinese glyphs (punctuation at the left of the em box), but the Traditional-
% Chinese fonts below centre ，：；。 etc., so compressed punctuation overlaps the
% next character. zhtw keeps those marks full width.
% ── CJK font: DFKai-SB (標楷體) on Windows, falling back to AR PL KaitiM Big5
% when DFKai-SB isn't installed (same setup as the midterm report).
\IfFontExistsTF{DFKai-SB}{%
  \setCJKmainfont{DFKai-SB}[JFM=zhtw]
  \setCJKsansfont{DFKai-SB}[JFM=zhtw]
  \setCJKmonofont{DFKai-SB}[JFM=zhtw]
}{%
  \setCJKmainfont{AR PL KaitiM Big5}[JFM=zhtw]
  \setCJKsansfont{AR PL KaitiM Big5}[JFM=zhtw]
  \setCJKmonofont{AR PL KaitiM Big5}[JFM=zhtw]
}

% ── Traditional-Chinese wording overrides ──────────────────────────────────────
\renewcommand{\figurename}{圖}
\renewcommand{\tablename}{表}
\renewcommand{\refname}{參考文獻}

% 版面（geometry、行距、頁碼）由各份文件自行設定。
% ══════════════════════════════════════════════════════════════════════════════

% ── JFM=zhtw (common/jfm-zhtw.lua): ctex's default JFM assumes simplified-
% Chinese glyphs (punctuation at the left of the em box), but the Traditional-
% Chinese fonts below centre ，：；。 etc., so compressed punctuation overlaps the
% next character. zhtw keeps those marks full width.
% ── CJK font: DFKai-SB (標楷體) on Windows, falling back to AR PL KaitiM Big5
% when DFKai-SB isn't installed (same setup as the midterm report).
\IfFontExistsTF{DFKai-SB}{%
  \setCJKmainfont{DFKai-SB}[JFM=zhtw]
  \setCJKsansfont{DFKai-SB}[JFM=zhtw]
  \setCJKmonofont{DFKai-SB}[JFM=zhtw]
}{%
  \setCJKmainfont{AR PL KaitiM Big5}[JFM=zhtw]
  \setCJKsansfont{AR PL KaitiM Big5}[JFM=zhtw]
  \setCJKmonofont{AR PL KaitiM Big5}[JFM=zhtw]
}

% ── Traditional-Chinese wording overrides ──────────────────────────────────────
\renewcommand{\figurename}{圖}
\renewcommand{\tablename}{表}
\renewcommand{\refname}{參考文獻}
